% AUTO-GENERATED -- do not edit by hand.
\begin{table}[H]
\centering
\small
\setlength{\tabcolsep}{3.2pt}
\begin{tabularx}{\linewidth}{llXrll}
\toprule
Specification / step & Universe / fixed axis & Horizon / comparison & $T$ & $\hat\rho$ [95\% CI] & $\hat\lambda$ [95\% CI] \\
\midrule
\multicolumn{6}{l}{\textit{Panel A: sequential specifications}} \\
\addlinespace[1.5pt]
Historical reference & 52 & Long (2023-01-03--2026-07-15) & 885 & 0.776 [0.743, 0.807] & 3.46 [2.89, 4.17] \\
Horizon-matched & 52 & Short (2023-01-03--2026-04-06) & 816 & 0.771 [0.735, 0.806] & 3.37 [2.77, 4.15] \\
Current governed & 100 & Short (2023-01-03--2026-04-06) & 815 & 0.843 [0.815, 0.869] & 5.36 [4.39, 6.61] \\
\midrule
\multicolumn{6}{l}{\textit{Panel B: descriptive step contrasts}} \\
\addlinespace[1.5pt]
Horizon & 52 firms & Long $-$ short & \textemdash & 0.005 & 0.091 \\
Universe & Short horizon & 100 $-$ 52 firms & \textemdash & 0.072 & 1.993 \\
\bottomrule
\end{tabularx}
\caption{Sequential universe-and-horizon bridge for the barycentric $W^\flat$ SAR QMLE. The historical first row uses the recovered 52-firm operator and long return window. The second row holds that operator and universe fixed while shortening the return endpoint to 6 April 2026; its interval is newly estimated with the joint-date stationary bootstrap. The third row expands to the governed current 100-firm operator at the same endpoint; the first and third intervals come from their verified governed results. Complete-case date counts can differ across universes. Panel B reports arithmetic differences only: they are descriptive, noncausal, and carry no contrast-level inference. $\hat\lambda=\hat\rho/(1-\hat\rho)$.}
\label{tab:p5-universe-horizon-bridge}
\end{table}
